\section{Adding an arm}
\label{sec:adding}

% Part of the contribution, not an appendix: a suite nobody can extend is not a
% suite. The two functions, the three conventions the shared code depends on,
% the commands that measure a new arm, and the baseline it compares against.
% The long form is docs/adding-an-arm.md in the repository; keep this section
% short enough that it cannot drift from that file.

A method is added by writing two functions in one file. The update kernel takes what
the method stores and the gradient the shared backward pass produced, and writes the
state and the next step's draw. The \texttt{uq\_sample\_fn} takes a sample index and
returns a device pointer to draw $s$. The rest is inherited.

Three conventions the shared code depends on, each of which fails silently. An arm's
\texttt{main} is compiled out when its file is included into the Python extension,
because four of them cannot coexist in one module, while its kernels and its callback
stay outside that guard so the module gets the arm's mathematics with no second copy.
cuBLAS is reached only through an accessor macro, which yields one handle per linked
image whatever the translation-unit layout. And a fatal path goes through
\texttt{UNCERTAIN\_ABORT}, \texttt{exit} in a binary and a thrown exception in the
module, so shared code cannot kill an interpreter.

A new arm is then measured by the commands the four here were: \texttt{bench} at one
batch size for a point on the sweep, \texttt{bench\_sweep.sh} for the whole of it, and
\texttt{eval} for the quality and uncertainty columns at a stated $S$ and split.
Tables~\ref{tab:posterior} and \ref{tab:quality} are the baseline it lands beside.
\texttt{docs/adding-an-arm.md} in the repository is the long form: which file to copy,
how the buffer layout is extended, how to count the passes a step makes, and what to
run before a number from a new arm is reported.
